\chapter{The Basics}
\newpage
\section{What a boomerang looks like}
Boomerang models can be classic (see figure \ref{fig:TechnicalTerms} and section \ref{sec:BasicShapes}), or unusual shapes (e.g. coathanger or scissor shape) to amaze people being introduced to boomerangs, or simply for the amusement of the thrower. Any boomerang may be appropriate for some type of competition. Boomerangs has many different shapes, still they have a few things in common.

\begin{figure}[ht]
	\centering
	\includegraphics[width=30em]{bterms.eps}
	\caption{Technical terms}
	\label{fig:TechnicalTerms}
\end{figure}

Each boomerang has a certain shape and set of airfoils that makes it fly the way it does (see figure \ref{fig:TechnicalTerms}). Left- and right-handers uses different boomerangs. As you may have noticed, Left-handed boomerangs are identical to the mirrored image of a right-handed.
\newpage

\section{Why does a boomerang come back}
I will only discuss the subject of physics briefly. If you would prefer a more detailed explanation, I recommend the work of Dr. Felix Hess, ``Boomerangs: Aerodynamics and Motion'' (see Literature list).
The following two questions are essentially for the return of a boomerang:
\begin{enumerate}
	\item Why is the flight path circular?
	\item Why does a boomerang thrown vertically lay-over to a horizontal position during flight?
\end{enumerate}
I will try to explain these phenomenons as simply as possible.

\paragraph{Why is the flight path circular?} 

First, we need to realize that the wings of a boomerang have airfoils that create ``lift'' (see figure \ref{fig:LiftAF})
\begin{figure}[h]
	\centering
	\includegraphics[width=20em]{04_LIFT.eps}
	\caption{Max lift on wing profile}
	\label{fig:LiftAF}
\end{figure}

\begin{center}
\colorbox{BOrange!30}{\parbox{30em}{
\colorbox{BOrange!30}{\parbox{28em}{
	\color{black}
	\vspace{1em}
	Now keep in mind the following two rules:
	\begin{itemize}
		\item[A)] The faster the air passes over the airfoils, the greater the lift.
		\item[B)] The better the angle of the air passing the wing (see figure \ref{fig:LiftBS}), the greater the lift.
	\end{itemize}
}}}}
\end{center}

\begin{figure}[h]
	\centering
	\includegraphics[width=20em]{05_LIFT.eps}
	\caption{Max lift on boomerang}
	\label{fig:LiftBS}
\end{figure}

To help us understand where a boomerang develops its lift let’s look at a construction disk (see figure \ref{fig:CD}). The circular lines represent air speed relative to the point around the boomerang is rotating at a moment in time. 
If you put the centre of rotation on the instant centre (i.e. forward speed equals zero) you can see how the air flows over the wings during the hover at the end of the flight. The faster the forward speed is in relation to the rotation speed (i.e. during the first part of the flight) the further the boomerang’s centre of rotation is from the instant centre. The centre of rotation is always on line V (see \ref{fig:BoomerangCD}).

The boomerang has a forward speed and rotation speed when thrown. Let’s look at an example to see how the forward and rotation speeds combine. We will assume that forward speed ($V_t$) equals the rotation speed ($V_r$). The boomerang is placed on the construction disk as shown in figure \ref{fig:SpeedRotComb}.
\begin{figure}[h]
	\centering
	\includegraphics[width=30em]{06_VTEVR.eps}
	\caption{Speed and rotation combined}
	\label{fig:SpeedRotComb}
\end{figure}
Let’s look at the total speed of any point on the boomerang above the centre of rotation at a moment in time. The forward speed is toward the left. Its rotation speed is also towards the left, so the total speed of the point is found by adding the two speeds \footnote{Actually these speeds need to be added as vectors with speed and direction (horizontal and vertical) considered. When this is done, for several points, the relative speeds shown by the concentric speed circles can be seen.}. When a wing is below the centre of rotation, its forward speed is still toward the left, but rotation speed is to the right, so the total speed is the forward speed minus the rotation speed. Because of the higher total speed of the wing above the centre of rotation, it generates more lift than the wing below the centre of rotation (see rule a). 

Since lift increases with speed, we would expect the maximum lift force in each wing to occur near the top of the diagram where the speed and spin are combined - near point $A$. Lift is greatest when the air travels at a right angle to the wing (see rule b). Points $A_1$ and $A_2$ represent points of maximum lift for the respective wings. If the boomerang did not behave as a gyroscope we would expect it to tilt toward the centre of the flight path. However, because of its rotation, it does behave as a gyroscope. Because of gyroscopic precession the lift force at $A$ has its effect 90° later at point $A^{\prime}$ (see figure \ref{fig:Gyro}a). This makes the boomerang turn left on its vertical axis, counter clockwise when viewed from above (see figure \ref{fig:Gyro}b).
\begin{figure}[h]
	\centering
		\subfigure[Effect of liftforce]{
			\includegraphics[height=10em]{07_VTEVR.eps}
		}
		\subfigure[Turning caused by gyroscopic precession]{
			\includegraphics[height=10em]{08_KREIS.eps}
		}
	\caption{Gyroscopic effect on Boomerang}
	\label{fig:Gyro}
\end{figure}
\newpage

\paragraph{Why does a boomerang lay-over?}
\begin{enumerate}
\item The lift of a wing occurs near the front - not in the centre of the airfoil (see figure \ref{fig:LiftAF}). Surprisingly, this is even true of symetrical airfoils (see figure \ref{fig:Lift}a). The point of average maximum lift $A_{\textrm{max}}$ is therefore moved slightly to the left (see figure \ref{fig:Lift}b).

\begin{figure}[h]
	\centering
	\subfigure[Lift on symetrical airfoil]{
	\includegraphics[height=5.5em]{09_LIFTS.eps}
	}
	\subfigure[$A_{\textrm{max}}$ is moved]{
	\includegraphics[height=5.5em]{10_LSHIF.eps}
	}
	\caption{Average lift}
	\label{fig:Lift}
\end{figure}

The actual turning force is therefore below the horizontal centreline at point $A^{\prime}$. The boomerang tends to roll clockwise when viewed from behind. This is called lay-over (see figure \ref{fig:Layover}). 
\begin{figure}[h]
	\centering
	\includegraphics[height=18em]{12_AKREIS.eps}
	\caption{Layover}
	\label{fig:Layover}
\end{figure}
\newpage

\item If the airfoil is symetrical the elbow normally causes more lift in front of the centre of rotation at Point E than it does during the other half of the rotation when it is behind the vertical centreline of the rotation disk (see figure \ref{fig:LiftElbow}).

\begin{figure}[h]
	\centering
	\includegraphics[width=30em]{11_ELLBO.eps}
	\caption{Lift acting on elbow}
	\label{fig:LiftElbow}
\end{figure}

 That is because of some shading effects to the inner elbow. The actual turning force, at Point $E^{\prime}$, is $90^{\circ}$ later, in the lower half of the rotation disk. When making two-bladers the airfoil of the elbow can be changed to alter the flight characteristics (for further information, see section \ref{sec:Elbow}). With symetrical multibladers the larger the area around the centre of rotation, the larger the lay-over.
\end{enumerate}
Because the resulting turning force is generally in the lower half of the rotation disk for most boomerangs, they lay-over in flight.
\newpage

\section{Making a boomerang} % Should be updated to include newer materials such as paxoline
First of all you need the appropriate material. One of the best materials is aircraft grade Finnish birch plywood. Typically, its plies are 0.5 mm thick. The digital boomerang plans refer to this material. Another really good material is paxoline. If you work in paxoline it is important you use dustmask and work in a well ventilated area.
Next you will need some tools: a jigsaw- or coping saw, a file, 80 and 180 grit sandpaper, 2 clamps, the edge of a table, a pencil, a paintbrush and paint. 

\begin{center}
\colorbox{BOrange!30}{\parbox{30em}{\centering
\colorbox{BOrange!30}{\parbox{28em}{
	\color{black}
\label{Sec:SafetyMaterials}
When using power tools such as a jigsaw or sander you should wear safety glasses. In addition you should wear a dust mask. When working with Fibre glass or Phenolic it is extremely important to use a GOOD filter mask, NOT just a dust mask!!
}}}}
\end{center}
Choose one of the provided plans. You can copy the shape and airfoil to make a boomerang. If you need to make a left-hander using a right-hander model, draw a mirror image of the outline and airfoil onto the plywood.

Follow the steps below to make your boomerang.
\begin{enumerate}
\item Choose the concave side of the plywood for the top side of the boomerang. This will ensure that the boomerang will not have negative dihedral (mentioned later in text). Draw the shape of the boomerang onto the plywood so that the grain of the wood runs from one wing to the other (see figure \ref{fig:PlyWoodGrains}). This adds strength to the elbow of the boomerang. 
\begin{figure}[h]
	\centering
	\includegraphics[width=30em]{13_MASER.eps}
	\caption{Wood Grains across boomerang}
	\label{fig:PlyWoodGrains}
\end{figure}
\item Saw out the shape of the boomerang with a jigsaw - or coping saw.
\item Draw the airfoil onto the piece of wood (see figure \ref{fig:AFonWood}).
\begin{figure}[h]
	\centering
	\includegraphics[width=30em]{14_BRECH.eps}
	\caption{Airfoil drawn onto the wood}
	\label{fig:AFonWood}
\end{figure}
\item Clamp the piece of wood onto the edge of a table.
\item Form your desired airfoil by filing. Ensure that the trailing edge is tapered and the leading edge is rounded (see figure \ref{fig:TechnicalTerms}).
\item Remove the file marks by sanding with 80 grit, then resand with 180 grit sandpaper.
\item Apply sanding sealer compatible with the finish of your choice and let it dry. 
\item Test throw your boomerang and make any changes you find necessary (see the chapter \ref{sec:DebugThrowing}). 
\item Resand with 180 grit sandpaper, then apply the decoration and finish coats. 
\end{enumerate}
\newpage

\section{Throwing}
\label{sec:Throwing}
The most important consideration is safety. Please read the following instructions carefuly.

\begin{center}
\colorbox{BOrange!30}{\parbox{30em}{
\begin{center}
\textbf{SAFETY}\\
\end{center}
\colorbox{BOrange!30}{\parbox{28em}{
	\color{black}
	\vspace{-1em}
\begin{enumerate}
 \item Length and width of your throwing area needs to be 2 to 3 times the range of the boomerang.
 \item You should never throw anywhere near people.
 \item If you are throwing in a group, only one boomerang at a time should be in the air.
 \item Make sure you warn people in time, if the boomerang flies towards them.
 \item Never take your eyes off a boomerang while it is flying.
\end{enumerate}
}}
}}
\end{center}
Most throwers throw from the middle of the field. Use a flag or throw a bit of grass into the air to check the wind direction and speed before each throw. If the wind is too strong (i.e. more than 4 Beaufort, table p. 43) go fly a kite instead.

\begin{figure}[h]
	\centering
	\subfigure[Pinch grip]{
	\includegraphics[height=14em]{16_GRIFF_B.eps}
	}
	\subfigure[Fore finger grip]{
	\includegraphics[height=14em]{16_GRIFF_A.eps}	
	}
	\caption{Common grips}	
	\label{fig:Grips}
\end{figure}

Now to the throw:
You decide which wing to throw from. Try different wings and do what works best for you.
The boomerang should be held between the index finger and thumb with the top of the boomerang toward you. There are two different grips, one is the pinch grip (see figure \ref{fig:Grips}a) and the other is the fore-finger grip (see figure \ref{fig:Grips}b). Try using both grips to see which works best for you.
When throwing, the following 3 angles are important.
\begin{enumerate}
\item The angle to the wind (see figure \ref{fig:CompassAngle})
\item The tilt angle (see figure \ref{fig:TiltAngle})
\item The height of throw (see figure \ref{fig:HorizonAngle})
\end{enumerate}

\begin{figure}[h]
	\centering
	\includegraphics[width=30em]{17_WINDR.eps}
	\caption{Angle to the wind (Compass angle)}	
	\label{fig:CompassAngle}
\end{figure}

It is important to give the boomerang enough spin when you throw. If you don’t it cannot return. Because all boomerangs fly differently, you may need to adjust one or more of your throwing angles to make a certain boomerang return (see chapter \ref{sec:DebugThrowing}.)
\newpage
\begin{figure}[p]
	\centering
	\includegraphics[height=20em]{18_TILT.eps}
	\caption{Tilt angle}	
	\label{fig:TiltAngle}
\end{figure}

\begin{figure}[p]
	\centering
	\includegraphics[width=30em]{19_HORI2.eps}
	\caption{The height of throw (angle above or below the horizon)}	
	\label{fig:HorizonAngle}
\end{figure}

\newpage
\section{Catching}
If your first attempts at throwing were successful, you can try to catch your boomerang. It should be lying flat toward the end of its flight and should approach you at a slow speed. If it was a particularly good throw, it almost appears as though the boomerang expects to be caught. The safest catch is a two-hand catch. Both palms should be held parallel to the boomerang as it hovers toward you. If they aren’t the boomerang may hit your fingers and spin back at you. Try to catch the boomerang by clapping the palms of your hands together, one hand below and the other above the boomerang (see figure \ref{fig:Catch}). 

\begin{figure}[h]
	\centering
	\includegraphics[height=20em]{20_CATCH.eps}
	\caption{Two methods of catching. Two- and one-handed.}	
	\label{fig:Catch}
\end{figure}

A one-hand catch is much more difficult. It is easiest when the boomerang hovers down to you. Two-arm boomerangs look as though there is a hole in the centre. You can either reach through the hole with your fingers or with your thumb to make the catch. Generally, multibladers are easier to catch one-handed. It is helpful to turn your hand in the direction of rotation as you make the catch (see figure \ref{fig:Catch}). I recommend that you wear gloves if you are new to catching boomerangs.
\newpage
\section{Wind}
Generally it is important to be careful when throwing in strong wind. A rule of thumb is 0-4 Beaufort is appropriate for most throwers. It is better to have a lot of stable wind rather than unpredictable and unsteady winds. 
\begin{center}
\colorbox{BOrange!30}{\parbox{30em}{
	\color{black}
	\vspace{1em}
	\centering
\begin{tabulary}{28em}{rLl}
Beaufort & Effect & m/s \\ [0.5ex] % inserts table \vspace{1em}
%heading
\hline                  % inserts single horizontal line
\\
\hspace*{-0.5 em}
0 & Calm & 0 - 0,2  \\ % inserting body of the table
1 & Slight draft, smoke does not rise completely upwards & 0,3 - 1,5  \\
2 & Slight breeze, noticeable & 1,6 - 3,3  \\
3 & Weak breeze, leaves flutter & 3,4 - 5,4 \\
4 & Steady breeze, twigs move & 5,5 - 7,9 \\
5 & Fresh breeze, movement in bigger twigs & 8,0 - 10,7\\
6 & Strong wind, movement in bigger branches & 10,8 - 13,8\\
7 & Stormy, movement in smaller trees trunks & 13,9 - 17,1\\
8 & Gale, movement in bigger trees trunks & 17,2 - 20,7\\
8 & Strong gale, roof tiles are lifted & 20,8 - 24,4\\
10 & Very strong gale, uproots trees & 24,5 - 28,4\\ [1ex]      % [1ex] adds vertical space
\end{tabulary}
}}
\end{center}
\newpage

\section{Accessories}
Apart from boomerangs, there are many useful accessories which should always be in your bag\\
\begin{center}
\colorbox{BOrange!30}{\parbox{30em}{
\colorbox{BOrange!30}{\parbox{28em}{
	\color{black}
	\vspace{1em}
\begin{enumerate}
\item Cycling gloves, goalkeeper gloves, wool gloves
	\begin{itemize}
		\item[-]	To protect your fingers while catching, and to keep your hands warm in cold weather
	\end{itemize}
\item Towels, cotton cloths, socks
	\begin{itemize}
		\item[-] For drying or cleaning the boomerangs
	\end{itemize}
\item Windflag, a stick with a piece of yarn attached
	\begin{itemize}
		\item[-] For reading the direction and speed of the wind
	\end{itemize}
\item Safety glasses / sunglasses
	\begin{itemize}
		\item[-] To protect your eyes from injury / excessive sunlight; also helps to see the boomerang in a bright sky
	\end{itemize}
\item Bandages
	\begin{itemize}
		\item[-] For small cuts; or to prevent wounds from throwing or catching
	\end{itemize}
\item Permanent marker (fine tip)
	\begin{itemize}
		\item[-] For making notes on the boomerang
	\end{itemize}
\item Stopwatch 
	\begin{itemize}
		\item[-] For timing
	\end{itemize}
\item Sandpaper, sanding block, file 
	\begin{itemize}
		\item[-] For modifying your boomerangs
	\end{itemize} 
\item Coins, pieces of lead and electrical tape
	\begin{itemize}
		\item[-] For adding weight to boomerangs 
	\end{itemize}
\end{enumerate}
}}
}}
\end{center}